% ==============================================================================
% CAPÍTULO 3 - METODOLOGIA
% ==============================================================================

\chapter{Metodologia}
\label{cap:metodologia}

O estágio foi realizado de forma presencial nas dependências da BIOTECNO, com carga horária distribuída
ao longo da semana conforme acordo com o supervisor Gustavo Henrique Buuron. As demandas de
desenvolvimento foram conduzidas de maneira iterativa e incremental, próxima a uma abordagem ágil
informal: as tarefas eram discutidas com o supervisor, executadas em ciclos curtos e revisadas à medida
que os resultados parciais emergiam, permitindo ajustes contínuos de escopo e implementação.

A organização das atividades seguiu um fluxo composto pelas seguintes etapas: (1) compreensão do
problema ou requisito apresentado pelo supervisor; (2) análise do sistema existente, leitura de
\textit{models}, \textit{views} e templates já implementados, para entender o contexto de integração;
(3) implementação da solução, com testes manuais no ambiente de desenvolvimento; e (4) revisão conjunta
com o supervisor antes de integrar a funcionalidade à base de código principal.

As ferramentas e tecnologias centrais do projeto foram:

\begin{alineas}
    \item \textbf{Python} e \textbf{Django}: linguagem e framework base do sistema web da empresa,
    utilizados em todas as etapas de desenvolvimento \cite{django2024};
    \item \textbf{LaTeX} com compilador \textbf{LuaLaTeX} via \texttt{latexmk}: sistema de composição
    tipográfica utilizado para a geração dos documentos PDF \cite{latex2023};
    \item \textbf{PostgreSQL}: sistema gerenciador de banco de dados relacional utilizado pela aplicação;
    \item \textbf{Git}: controle de versão para rastreamento das alterações no código-fonte;
    \item \textbf{Django Admin}: interface administrativa nativa do Django, utilizada como ponto de
    gerenciamento dos layouts de orçamento.
\end{alineas}

A comunicação técnica com o supervisor ocorreu diariamente, permitindo que decisões de arquitetura 
como a escolha de resolver o layout ativo em tempo de geração em vez de vinculá-lo ao modelo
\texttt{Orcamento} fossem validadas antes de serem implementadas. Essa dinâmica reduziu retrabalho e
garantiu que a solução fosse coerente com as necessidades operacionais da empresa.